\textbf{Equal information.} The comparison above favours CNA+, which calibrates on all training units while the uncertified policies were tuned on the 60/40 split. The third protocol therefore rebuilt the comparators on exactly the same five fold models: selected on the out-of-fold predictions of all training units and applied through the average of the five models (CF in Table~\ref{tab:cost}). With equal information the certified alarm cost about as much as the best uncertified rules (H9). Its mean excess over the cheapest of the six policies in this comparison was 1.0--14.3\,\% across the seven datasets, against 0.6--11.0\,\% for the one-standard-error cost-tuned alarm and 1.8--13.5\,\% for policy~1 of Kamariotis et al.\ with an optimised threshold. Averaged over the three repetitions (descriptive) the three ranges were 1.9--14.6, 0.9--11.7 and 1.9--11.8\,\%, and the excess of CNA+-DA moved by up to 12 points between repetitions (N-CMAPSS), so differences of a few points among these three policies are within the variability of the split. By the protocol's cell-count rule CNA+-DA was never dearer than either: it was cheaper than the one-standard-error alarm on the battery cells (16 of 16 cells) and on FD004 (8 of 16) and mixed on the other five datasets, and mixed against the optimised probability-threshold policy on all seven. It was cheaper than policy~1 with the heuristic threshold $1/\rho$ on PHM08, FD002 and FD004. Plain cost tuning won many cells by small margins but lost heavily in a few, because it occasionally chose a setting at the edge of the feasible region: its largest excess in a cell was 97\,\% on FD003, against 45\,\% for CNA+-DA. Where CNA+-DA did lose, it lost at $\rho=50$: with about 80 calibration units its failure probability cannot fall below about 1.2\,\% (Proposition~\ref{prop:cost}), and the one-standard-error alarm, which can choose levels beyond the most demanding calibration unit, had no test failure in those cells. The one-standard-error and optimised-threshold comparators failed on 0--1\,\% of test units at the primary lead time and $\rho=10$ (all three repetitions), the heuristic threshold on 3.3--13.1\,\%; the former achieved low failure rates on these units, but without a guarantee for the next one.
